\documentclass[12pt, a4paper]{article}
\usepackage{../../../myconfig} % Gọi file cấu hình từ ngoài gốc vào

% ==========================================
% BẮT ĐẦU NỘI DUNG TÀI LIỆU
% ==========================================
\begin{document}

% Truyền 3 tham số: {Nhãn cam} {Chữ to chính giữa} {Chữ ở góc phải Header}
\titlebox{Chủ đề: Hàm số}{Tốc độ thay đổi của một đại lượng}{Hàm số: Tốc độ thay đổi đại lượng}

% --- CÂU 6 ---
\questionTrueFalse{0}{1}{Người ta mô hình hóa nhiệt độ trung bình hàng tháng ở thành phố Honolulu, Hawaii bằng hàm số \( F(t) = -0,0418t^3 + 0,605t^2 - 1,16t + 73,4 \) (đơn vị: độ Fahrenheit), với \( t \) là số tháng, \( t = 0 \) ứng với thời điểm đầu tháng 1/2024. Cho biết công thức đổi từ độ Fahrenheit sang độ Celsius là \( F = \dfrac{9}{5}C + 32 \), với \( F \) và \( C \) lần lượt là độ Fahrenheit và độ Celsius.}{
    \textbf{a)} & Trong năm 2024, nhiệt độ trung bình hàng tháng ở thành phố Honolulu tăng liên tục trong hơn 7 tháng liên tiếp.
    & \textcolor{amberorange}{$\bigcirc$} & \textcolor{amberorange}{$\bigcirc$} \\ \hline
    
    \textbf{b)} & Nhiệt độ trung bình cao nhất trong năm 2024 ở thành phố Honolulu là \( 27,5^\circ C \) (Làm tròn kết quả đến hàng phần mười của độ Celsius).
    & \textcolor{amberorange}{$\bigcirc$} & \textcolor{amberorange}{$\bigcirc$} \\ \hline
    
    \textbf{c)} & Vào cuối tháng 8/2024, nhiệt độ trung bình hàng tháng ở thành phố Honolulu đang giảm với tốc độ là \( 0,49^\circ F/\text{tháng} \) (kết quả làm tròn đến hàng phần trăm).
    & \textcolor{amberorange}{$\bigcirc$} & \textcolor{amberorange}{$\bigcirc$} \\ \hline
    
    \textbf{d)} & Nhiệt độ trung bình hàng tháng ở thành phố Honolulu tăng nhanh nhất vào khoảng cuối tháng 5/2024.
    & \textcolor{amberorange}{$\bigcirc$} & \textcolor{amberorange}{$\bigcirc$} \\
}

% --- CÂU 7 ---
\questionTrueFalse{0}{2}{Để khắc phục tình trạng thất nghiệp tại quốc gia X, chính phủ của quốc gia này quyết định sẽ đưa ra các gói trợ cấp để tuyển dụng những người đang thất nghiệp vào làm việc cho các dự án của chính phủ.  
Giả sử rằng sau \( t \) tháng kể từ khi triển khai việc cung cấp các gói trợ cấp, khi đó quốc gia X có  
\(N(t) = -t^3 + 45t^2 + 408t + 3078\) 
người thất nghiệp. Cho biết rằng việc cung cấp các gói trợ cấp này được triển khai trong 4 năm.}{
    \textbf{a)} & Sau 3 năm kể từ khi triển khai việc cung cấp các gói trợ cấp, quốc gia X có 4680 người thất nghiệp.
    & \textcolor{amberorange}{$\bigcirc$} & \textcolor{amberorange}{$\bigcirc$} \\ \hline
    
    \textbf{b)} & Trong thời gian triển khai việc cung cấp các gói trợ cấp, số người thất nghiệp nhiều nhất tại cùng một thời điểm là 29699 người.
    & \textcolor{amberorange}{$\bigcirc$} & \textcolor{amberorange}{$\bigcirc$} \\ \hline
    
    \textbf{c)} & Sau 1 năm kể từ khi triển khai việc cung cấp các gói trợ cấp, tốc độ tăng số người thất nghiệp tại thời điểm đó là 1056 người/tháng.
    & \textcolor{amberorange}{$\bigcirc$} & \textcolor{amberorange}{$\bigcirc$} \\ \hline
    
    \textbf{d)} & Số người thất nghiệp tại quốc gia X tăng nhanh nhất sau 1 năm 3 tháng kể từ khi triển khai việc cung cấp các gói trợ cấp.
    & \textcolor{amberorange}{$\bigcirc$} & \textcolor{amberorange}{$\bigcirc$} \\
}

% --- CÂU 8 ---
\questionTrueFalse{0}{3}{Tin đồn là những thông tin chưa được kiểm chứng, không có nguồn gốc xác thực nhưng vẫn đạt được trạng thái "viral". Tin đồn càng ảnh hưởng, nằm trong phạm vi quan tâm của nhiều người và càng khó kiểm chứng thì tính hấp dẫn với cộng đồng càng cao, lôi kéo được nhiều người và tăng được tính lan tỏa.  
Giả sử tỉ lệ phần trăm học sinh trong một trường học biết về một tin đồn mới được lan truyền được mô hình hóa bởi hàm số  
\(p(t) = \dfrac{100}{1 + 5e^{-0,2t}}\)
trong đó \( t \) là số giờ kể từ khi tin đồn bắt đầu được lan truyền. Biết rằng tin đồn này được bắt đầu lan truyền vào lúc 7 giờ ngày 26/06/2026, và ngôi trường này có 2000 học sinh.}{
    \textbf{a)} & Vào lúc 10 giờ ngày 26/06/2026, có khoảng 543 học sinh đã biết đến tin đồn (kết quả làm tròn đến hàng đơn vị).
    & \textcolor{amberorange}{$\bigcirc$} & \textcolor{amberorange}{$\bigcirc$} \\ \hline
    
    \textbf{b)} & Vào lúc 10 giờ ngày 26/06/2026, tốc độ lan truyền của tin đồn này là \( 3,92\% \) (kết quả làm tròn đến hàng phần trăm).
    & \textcolor{amberorange}{$\bigcirc$} & \textcolor{amberorange}{$\bigcirc$} \\ \hline
    
    \textbf{c)} & Tốc độ lan truyền của tin đồn này đạt giá trị lớn nhất sau \( a \ln b \) giờ kể từ khi tin đồn bắt đầu được lan truyền (\( a; b \in \mathbb{Z}, b \) là số nguyên tố). Khi đó, giá trị của \( a + b = 10 \).
    & \textcolor{amberorange}{$\bigcirc$} & \textcolor{amberorange}{$\bigcirc$} \\ \hline
    
    \textbf{d)} & Vào thời điểm tốc độ lan truyền của tin đồn này đạt giá trị lớn nhất, khi đó tin đồn này đã có 1000 học sinh biết đến.
    & \textcolor{amberorange}{$\bigcirc$} & \textcolor{amberorange}{$\bigcirc$} \\
}


% --- CÂU 16 ---
\questionShortAnswer{0}{4}{Một bao cát nặng 100g bị rách ở đáy bao. Người ta mô hình hóa lượng cát còn lại trong bao cát này sau \( t \) giây kể từ khi bao cát bị rách bởi hàm số  
\(S(t) = 100\left(1 - \dfrac{t^2}{15}\right)^3\)
(đơn vị: gam). Biết rằng cát chảy ra khỏi bao nhanh nhất tại thời điểm \( t = a \) giây kể từ khi bao cát bị rách. Giá trị của \( a^6 \) bằng bao nhiêu?}

% --- CÂU 17 ---
\questionShortAnswer{0}{5}{Một người quản lý ở một trại nuôi cá xác định rằng sau \( t \) tuần kể từ khi thả 500 con cá của một chúng cá vào một chiếc hồ, cân nặng trung bình của một con cá (đơn vị: kg) được cho bởi công thức  
\(w(t) = -0,025t^2 + 0,6t + 1,5.\)
Ngoài ra, người này còn thấy rằng tỉ lệ cá còn sống sót sau \( t \) tuần được cho bởi công thức  
\(p(t) = \dfrac{25}{25 + t}.\)
Cho biết tổng sản lượng của loại cá này được tính bằng tổng khối lượng của những con cá còn sống trong ao tại thời điểm \( t \). Hỏi tổng sản lượng của loại cá này đạt giá trị lớn nhất là bao nhiêu tấn (kết quả làm tròn đến hàng phần trăm)?}

% --- CÂU 18 ---
\questionShortAnswer{0}{6}{Một cuộc khảo sát về tình hình ủng hộ của nhà chính trị gia MIN cho thấy rằng sau khi tham gia một cuộc tranh luận với nhà chính trị gia MAX, tỉ lệ phần trăm dân số ủng hộ cho MIN sẽ thay đổi theo hàm số  

\(S(t) = \dfrac{100(t^2 - 3t + 25)}{t^2 + 7t + 25},\)

trong đó \( t \) là số ngày tính từ khi cuộc tranh luận diễn ra. Biết rằng cuộc bầu cử chính thức sẽ được diễn ra vào thời điểm \( t = 10 \). Hỏi tỉ lệ phần trăm dân số ủng hộ cho MIN tăng nhanh nhất sau bao nhiêu ngày kể từ khi cuộc tranh luận diễn ra (kết quả làm tròn đến hàng phần trăm)?}

% --- CÂU 19 ---
\questionShortAnswer{0}{7}{Nhiệt độ cơ thể \( T \) (tính bằng độ Celsius) của một người bị cảm có thể được mô hình hóa bởi hàm số  
\(T(t) = \dfrac{5}{9} \left( \dfrac{12t}{t^2 + 8} + 66,6 \right),\)
trong đó \( t \) là số giờ kể từ khi một người bắt đầu có dấu hiệu của việc bị cảm.  
Biết rằng tại thời điểm \( t = a \) giờ thì nhiệt độ cơ thể của người này là cao nhất, và tại thời điểm \( t = b \) giờ thì nhiệt độ cơ thể của người này giảm nhanh nhất. Giá trị của \( a + b \) bằng bao nhiêu (kết quả làm tròn đến hàng phần trăm)?}



% --- CÂU 22 ---
\questionShortAnswer{0}{8}{Năm 1885, nhà tâm lý học Hermann Ebbinghaus đưa ra lý thuyết về đường cong quên lãng. Theo đó, khi học tập, mọi thông tin mới tiếp xúc đều sẽ được lưu vào khu trí nhớ ngắn hạn. Nếu như không được nhắc lại hay sử dụng thường xuyên, các thông tin này sẽ bị lãng quên nhanh; từ đó hình thành đường cong quên lãng. Cho biết đường cong quên lãng được mô hình hóa cho con người có dạng  
\(p(t) = (100 - a)e^{-bt} + a,\)
trong đó \( p(t) \) là tỉ lệ phần trăm thông tin mà một người còn nhớ được sau \( t \) tuần; còn \( a \) và \( b \) là các tham số ứng với từng người.  

Giả sử rằng sau khi một học sinh của anh Huy Math-ster học xong một chương kiến thức mới và không ôn tập lại, sau 1 tuần học sinh đó đã quên mất 47,5\% nội dung của chương, và tỉ lệ này sau 2 tuần tăng lên thành 71,25\%.  
Khi đó, tốc độ thay đổi lượng thông tin nhớ được của học sinh đó sau 3 tuần kể từ khi học xong chương kiến thức là bao nhiêu phần trăm/tuần (kết quả làm tròn đến hàng phần mười)?}


% --- CÂU 24 ---
\questionShortAnswer{0}{9}{Nhiệt độ của một nồi trà sau khi bỏ vào tủ đông mềm \( t \) phút được mô hình hóa bởi hàm số  
\(T(t) = \dfrac{1300}{t^2 + 2t + 25}\)
(đơn vị: độ Celsius). Hỏi sau bao nhiêu phút kể từ khi nồi trà được bỏ vào tủ đông mềm thì nhiệt độ của nó sẽ giảm nhanh nhất (kết quả làm tròn đến hàng phần trăm)?}

% --- CÂU 25 ---
\questionShortAnswer{0}{10}{Giả sử doanh số bán hàng (đơn vị triệu đồng) của một sản phẩm điện tử mới được mô hình hóa bằng hàm số  
\( f(t) = 600 \left( t^2 + ne^{-t} \right), \) 
với \( t \geq 0 \) là thời gian tính bằng năm kể từ khi phát hành sản phẩm mới, và \( n \leq 0 \) là tham số. Khi đó đạo hàm \( f'(t) \) biểu thị tốc độ bán hàng. Biết rằng tốc độ bán hàng luôn tăng trong khoảng thời gian 8 năm đầu từ khi ra mắt sản phẩm. Hỏi giá trị nhỏ nhất của \( n \) là bao nhiêu để điều kiện này được thỏa mãn?}

\questionShortAnswer{20}{11}{Sự tăng trưởng của một loại vi khuẩn trong môi trường nuôi cấy sẽ tuân theo công thức \( N = N_0 e^{rt} \), trong đó \( N_0 \) là số lượng vi khuẩn ban đầu, \( r \) là tỉ lệ tăng trưởng (\( r > 0 \)), \( t \) (giờ) là thời gian tăng trưởng. Biết rằng số lượng vi khuẩn ban đầu là 500 con và sau 3 giờ thì có 3000 con vi khuẩn. Hỏi sau bao nhiêu giờ thì tốc độ tăng số lượng vi khuẩn là 1000 con/giờ (kết quả làm tròn đến hàng phần trăm)?}

% --- CÂU 13 ---
\questionShortAnswer{20}{12}{Tỷ lệ trao đổi chất trong cơ thể của một người vừa hoàn thành một bữa ăn sẽ tăng dần, sau đó sẽ giảm dần sau một thời gian để cơ thể trở về trạng thái nghỉ. Hiện tượng này được gọi là hiệu ứng nhiệt của thực phẩm. Các nhà nghiên cứu đã chỉ ra rằng hiệu ứng nhiệt của thực phẩm cho một người trưởng thành được tính bởi công thức \( F(t) = -10,28 + 175,9t e^{-1,3t} \), với \( F(t) \) là hiệu ứng nhiệt của thực phẩm (đơn vị: kJ/h) và \( t \) (đơn vị: giờ, \( t > 0 \)) là thời gian trôi qua tính từ khi hoàn thành bữa ăn. Hỏi tốc độ thay đổi hiệu ứng nhiệt của thực phẩm đạt giá trị nhỏ nhất sau bao nhiêu phút kể từ khi hoàn thành bữa ăn?}

% --- CÂU 14 ---
\questionShortAnswer{20}{13}{Một nhóm các nhà nghiên cứu về xu hướng hành vi của con người nhận thấy rằng với những đoạn video mang tính chất học thuật, người xem sẽ thường có xu hướng xem những video có thời lượng vừa phải, do những video quá ngắn thì thường không cung cấp đủ thông tin, còn những video quá dài thì thường sẽ dễ gây ra sự nhàm chán. Với một video nói về toán học, các nhà nghiên cứu nhận thấy rằng tỉ suất người xem của một video kéo dài \( t \) phút được cho bởi hàm số  
\(R(t) = \dfrac{25t}{t^2 + 100}\) 
(đơn vị: phần trăm, \( t > 0 \)). Hỏi với độ dài của video toán học là bao nhiêu phút thì tỉ suất người xem video sẽ giảm nhanh nhất (kết quả làm tròn đến hàng phần mười)?}

% --- CÂU 15 ---
\questionShortAnswer{20}{14}{Mối liên hệ giữa cân nặng của một con nai sừng tấm cái với độ tuổi của nó có thể được biểu diễn bởi hàm số  
\(M(t) = 369 \cdot 0,93^t \cdot t^{0,36} \quad (0 \leq t \leq 15),\)  
với \( M(t) \) là cân nặng của một con nai sừng tấm cái (đơn vị: kg) và \( t \) là độ tuổi của con nai sừng tấm (tính bằng năm). Hỏi sau bao nhiêu năm kể từ khi được sinh ra thì con nai sừng tấm cái sẽ có tốc độ thay đổi cân nặng nhanh nhất (kết quả làm tròn đến hàng phần mười)?}



\end{document}